% Figure: stress concentration at a circular hole in a plate under remote
% uniaxial tension. The normalized stress sigma_yy / sigma_0 along the x-axis
% from the hole edge (Kirsch solution): at edge (r=a) it is 3, decaying to 1.
\begin{figure}[htbp]
\centering
\begin{tikzpicture}
\begin{axis}[
  width=12cm, height=7.5cm,
  xlabel={distance from hole centre, $x/a$},
  ylabel={normalised stress $\sigma_{yy}/\sigma_0$},
  xmin=1, xmax=5, ymin=0, ymax=3.4,
  axis lines=left,
  tick label style={font=\scriptsize},
  label style={font=\small},
  legend pos=north east,
  legend style={font=\scriptsize},
]
% Kirsch along the line perpendicular to load (theta=90deg), on the x-axis:
% sigma_yy/sigma_0 = 1 + 0.5 (a/x)^2 + 1.5 (a/x)^4
\addplot[engblue, very thick, samples=100, domain=1:5]
  {1 + 0.5*(1/x)^2 + 1.5*(1/x)^4};
\addlegendentry{Kirsch solution}
% asymptote at 1
\addplot[enggray, dashed, samples=2, domain=1:5] {1};
\addlegendentry{remote stress $\sigma_0$}
% mark Kt=3 at edge
\filldraw[engred] (axis cs:1,3) circle [radius=1.8pt];
\node[engred, font=\scriptsize, anchor=west] at (axis cs:1.1,3) {$K_t=3$ at edge};
\end{axis}
\end{tikzpicture}
\caption[Stress concentration at a circular hole]{Stress along the
line perpendicular to the load, outward from the edge of a circular
hole of radius $a$ in a wide plate under remote uniaxial tension
$\sigma_0$ (Kirsch solution, 1898). At the hole edge ($x=a$) the
stress is $3\sigma_0$, so the stress concentration factor is $K_t=3$,
independent of hole size. The amplification decays rapidly with
distance, falling to within ten percent of the remote stress by about
three radii out, an instance of Saint-Venant's principle. The factor
of three is the canonical number every mechanical designer carries:
a small hole triples the local stress.}
\label{fig:vol03ch03:stress-concentration-hole}
\end{figure}
